\section{Тест производительности}

Для оценки производительности реализован отдельный бенчмарк, сравнивающий словарь на собственном красно-чёрном дереве со стандартным \texttt{std::map}. В измерение входят только операции над структурами данных \texttt{Insert}, \texttt{Find} и \texttt{Erase}; загрузка и сохранение в файл не тестируются, так как они зависят уже не только от структуры словаря, но и от работы файловой системы.

Тестовые данные генерируются скриптом \texttt{generator.py}. Для каждого запуска сначала выполняется \texttt{prefill = max(5000, n / 10)} стартовых вставок без измерения времени, после чего измеряется серия из $n$ операций. Для каждого размера и сценария выполнено 5 прогонов, в таблицах приведены медианные значения времени в микросекундах. В последнем столбце указано отношение \texttt{std::map / RB}; значение больше 1 означает, что собственная реализация быстрее.

Рассматривались три сценария нагрузки:
\begin{itemize}
\item \texttt{mixed}: 50\% поисков, 25\% вставок, 25\% удалений;
\item \texttt{read-heavy}: 80\% поисков, 10\% вставок, 10\% удалений;
\item \texttt{update-heavy}: 40\% вставок, 40\% удалений, 20\% поисков.
\end{itemize}

Команда запуска полной серии измерений:
\begin{alltt}
$ make benchmark
$ bash run_benchmarks.sh benchmark_results.csv
\end{alltt}

\textbf{Смешанный сценарий:}
\begin{longtable}{|r|r|r|r|}
\hline
\rowcolor{lightgray}
$n$ & RB-tree, \texttt{us} & \texttt{std::map}, \texttt{us} & \texttt{std::map / RB} \\
\hline
50\,000 & 20\,026 & 20\,811 & $1.039\times$ \\
\hline
100\,000 & 88\,916 & 153\,207 & $1.723\times$ \\
\hline
200\,000 & 185\,440 & 211\,016 & $1.138\times$ \\
\hline
500\,000 & 292\,434 & 332\,745 & $1.138\times$ \\
\hline
1\,000\,000 & 1\,327\,619 & 1\,292\,859 & $0.974\times$ \\
\hline
\end{longtable}

В смешанном сценарии собственная реализация оказалась быстрее в 4 случаях из 5. Наиболее заметный выигрыш получился при $n = 100\,000$, где медианное время оказалось лучше примерно в $1.72\times$. На максимальном размере входа результаты почти сравнялись, и небольшое преимущество осталось уже у \texttt{std::map}.

\textbf{Сценарий с преобладанием поиска:}
\begin{longtable}{|r|r|r|r|}
\hline
\rowcolor{lightgray}
$n$ & RB-tree, \texttt{us} & \texttt{std::map}, \texttt{us} & \texttt{std::map / RB} \\
\hline
50\,000 & 20\,361 & 19\,646 & $0.965\times$ \\
\hline
100\,000 & 46\,539 & 48\,844 & $1.050\times$ \\
\hline
200\,000 & 86\,936 & 91\,528 & $1.053\times$ \\
\hline
500\,000 & 437\,522 & 487\,610 & $1.114\times$ \\
\hline
1\,000\,000 & 1\,083\,802 & 1\,110\,723 & $1.025\times$ \\
\hline
\end{longtable}

При преобладании операций поиска обе реализации показывают очень близкие результаты. Начиная с $10^5$ операций собственное дерево немного опережает \texttt{std::map}, но преимущество здесь уже не столь велико: в основном речь идёт о нескольких процентах.

\textbf{Сценарий с преобладанием модификаций:}
\begin{longtable}{|r|r|r|r|}
\hline
\rowcolor{lightgray}
$n$ & RB-tree, \texttt{us} & \texttt{std::map}, \texttt{us} & \texttt{std::map / RB} \\
\hline
50\,000 & 22\,733 & 24\,350 & $1.071\times$ \\
\hline
100\,000 & 49\,666 & 49\,653 & $1.000\times$ \\
\hline
200\,000 & 187\,465 & 145\,675 & $0.777\times$ \\
\hline
500\,000 & 417\,693 & 387\,555 & $0.928\times$ \\
\hline
1\,000\,000 & 2\,104\,113 & 1\,952\,622 & $0.928\times$ \\
\hline
\end{longtable}

Наиболее трудным для собственной реализации оказался сценарий, в котором доминируют вставки и удаления. На малых входах различие почти отсутствует, но начиная с $n = 200\,000$ стандартная библиотечная реализация работает быстрее. Это показывает, что одинаковая асимптотика $O(\log n)$ ещё не гарантирует одинаковые константы: при большом числе модифицирующих операций качество конкретной реализации становится особенно заметным.

В целом, результаты можно считать хорошими: собственное красно-чёрное дерево оказалось конкурентоспособным по сравнению с \texttt{std::map} и в двух сценариях из трёх чаще показывало равное или лучшее время. Одновременно бенчмарк показал и ограничение реализованного решения: под интенсивной нагрузкой на вставку и удаление стандартная библиотека обычно выигрывает за счёт более отлаженной и оптимизированной реализации.
\pagebreak
